% =====================================================================
%  IMPA Statistics & AI seminar — 2026-05-18 (Monday, 13:30)
%  "Log log plot — Estimating Critical Exponents"
%
%  Structural skeleton — one frame stub per planned slide.
%  Every frame currently holds a red \TODO{...} placeholder; content
%  lands in a later pass once the §-by-§ discussions happen.
%
%  Build:  latexmk -pdf talk.tex
% =====================================================================
\documentclass[10pt, aspectratio=169, xcolor={table}]{beamer}

% ---------- theme ----------------------------------------------------
%  Boadilla: white background, small footer with author/title/page —
%  the standard IMPA look. The colour scheme below paints all
%  structural details in dark blue (IMPA accent).
\usetheme{Boadilla}

% ---- IMPA colour scheme --------------------------------------------
% Invisible / Medium / Dark variants of the three traffic-light hues,
% then bound to Beamer's structural colour slots.
\definecolor{InvisibleRed}{rgb}{0.97, 0.92, 0.92}
\definecolor{InvisibleGreen}{rgb}{0.92, 0.97, 0.92}
\definecolor{InvisibleBlue}{rgb}{0.92, 0.92, 0.97}

\definecolor{MediumRed}{rgb}{0.925, 0.345, 0.345}
\definecolor{MediumGreen}{rgb}{0.37, 0.7, 0.66}
\definecolor{MediumBlue}{rgb}{0.015, 0.315, 0.45}

\definecolor{DarkBlue}{rgb}{0.05, 0.15, 0.35}

\usecolortheme[named=DarkBlue]{structure}

\setbeamercolor{titlelike}{parent=structure}
\setbeamercolor{block title}{bg=MediumBlue, fg=white}
\setbeamercolor{block body}{bg=InvisibleBlue}
\setbeamercolor{block title example}{bg=MediumGreen, fg=white}
\setbeamercolor{block body example}{bg=InvisibleGreen}
\setbeamercolor{block title alerted}{bg=MediumRed, fg=white}
\setbeamercolor{block body alerted}{bg=InvisibleRed}

% ---------- packages -------------------------------------------------
\usepackage[utf8]{inputenc}
\usepackage[T1]{fontenc}
\usepackage{amsmath, amssymb, amsthm, mathtools}
\usepackage{graphicx}
\usepackage{xcolor}
\usepackage{tikz}
\usepackage{booktabs}
\usepackage[backend=biber, style=numeric]{biblatex}
\addbibresource{refs.bib}

\graphicspath{{../images/}}

% ---------- footer override -----------------------------------------
% Slim two-column footer: "log-log plot" on the left, "xx/38" on the
% right. Overrides the default Boadilla three-segment footer.
\setbeamertemplate{navigation symbols}{}
\setbeamertemplate{footline}{%
  \leavevmode%
  \hbox{%
    \begin{beamercolorbox}[wd=0.5\paperwidth, ht=2.25ex, dp=1ex,
                           leftskip=1em]{author in head/foot}%
      \usebeamerfont{author in head/foot}log-log plot%
    \end{beamercolorbox}%
    \begin{beamercolorbox}[wd=0.5\paperwidth, ht=2.25ex, dp=1ex,
                           rightskip=1em]{date in head/foot}%
      \hfill\usebeamerfont{date in head/foot}%
        \insertframenumber/\inserttotalframenumber%
    \end{beamercolorbox}%
  }%
  \vskip0pt%
}

% ---------- macros (minimal) -----------------------------------------
\newcommand{\E}{\mathbb{E}}
\renewcommand{\P}{\mathbb{P}}
\newcommand{\Var}{\operatorname{Var}}

% Notes-to-self — visible while drafting; will be stripped before the talk.
\newcommand{\TODO}[1]{{\color{red}\textbf{[TODO]} #1}}

% Advisor / co-advisor (custom — Beamer has no built-ins for these).
\newcommand{\theadvisor}{}
\newcommand{\thecoadvisor}{}
\newcommand{\advisor}[1]{\renewcommand{\theadvisor}{#1}}
\newcommand{\coadvisor}[1]{\renewcommand{\thecoadvisor}{#1}}

% ---------- demo-break icon (TikZ) -----------------------------------
% Stylised play button inside a rounded "screen", with a label below.
% Used for the GIF / live-demo transition slides.
\newcommand{\demoicon}[1]{%
  \begin{center}
    \begin{tikzpicture}
      \draw[line width=3pt, rounded corners=10pt,
            fill=InvisibleBlue, draw=DarkBlue]
        (-3,-1.8) rectangle (3,1.8);
      \fill[MediumBlue] (-0.85,-1.05) -- (-0.85,1.05) -- (1.15,0) -- cycle;
      \node[below=0.45cm, font=\Large\bfseries, text=DarkBlue]
        at (0,-1.8) {#1};
    \end{tikzpicture}
  \end{center}}

\newcommand{\demoframe}[2]{%
  \begin{frame}[plain]
    \vfill
    \demoicon{#1}
    \vspace{1em}
    \begin{center}\large #2\end{center}
    \vfill
  \end{frame}}

% =====================================================================
%  Title block
% =====================================================================
\advisor{Augusto Teixeira}
\coadvisor{Roberto Imbuzeiro}

\title{Log log plot:}
\subtitle{Estimating Critical Exponents}
\author{Barbosa I.\\[0.6em]
  {\footnotesize Advisor: \theadvisor\\
   Co-advisor: \thecoadvisor}}
\institute{IMPA -- Instituto Nacional de Matemática Pura e Aplicada}
\date{18 May 2026}
\titlegraphic{\hfill\includegraphics[height=1.4cm]{impa-logo.png}\hspace{0.4em}}

% =====================================================================
\begin{document}

\maketitle

\begin{frame}{Plan of the talk}
  \tableofcontents
\end{frame}

% =====================================================================
\section{Phase transitions in percolation}
% target: ≈ 7 min
% =====================================================================

\begin{frame}[plain]
  \begin{center}
    \includegraphics[height=0.95\paperheight, keepaspectratio]%
                    {Broadbent_Hammersley\(1957\).pdf}
  \end{center}
  \note{%
    Broadbent \& Hammersley (1957)~\cite{broadbentHammersley1957}
    wanted to model a gas mask: a porous carbon granule with random
    defects. They threw away \emph{everything} physical except one
    number --- the density $p$ of \emph{open} sites.\par\smallskip
    %
    No particle interactions. No movement. No gravity. No chemistry.
    Just: each site is open with probability $p$, independently; a
    ``fluid'' percolates iff the open cluster of the source reaches
    the surface.\par\smallskip
    %
    Caricature of a physical system $\Rightarrow$ surprisingly rich
    mathematics.%
  }
\end{frame}

\demoframe{SIMULATION}{run - \texttt{perc\_interactive.py}}

\begin{frame}{Phase transitions everywhere}
  \begin{itemize}
    \item Ising model (magnetisation)
    \item XY model
    \item Diffusion / interacting particle systems
    \item Random walks
    \item Spread of infection
    \item Polymers / self-avoiding walks
  \end{itemize}
  \vspace{0.6em}
  Common thread: a control parameter, an order parameter, a critical
  value, power-law behaviour at the threshold.
  \vspace{0.6em}

  \TODO{Verbal: \emph{the same drama plays out in many
        statistical-mechanics models. We restrict ourselves to
        percolation today, but the wider problem is to do statistics
        on any of these.}
        Drop in image / icon thumbnails next to each bullet later.}
\end{frame}

\begin{frame}{Simulating percolation on a computer}
  \TODO{Chalkboard explanation --- slide is a stub.}

  \vspace{0.6em}
  \textbf{Uniform coupling.} Draw one i.i.d.\ array
  $U_x \sim \mathrm{Unif}[0,1]$ once; for any $p$, declare site $x$
  open iff $U_x \le p$. Same realisation, every $p$ simultaneously
  $\Rightarrow$ monotone in $p$.

  \TODO{Perfect for the slider demo.

        \textbf{(2) 3/4/6-connected.} Same recipe, different
        neighbourhoods: honeycomb ($z=3$), square ($z=4$), triangular
        ($z=6$). Coordination number $z$ controls how easy it is for
        the cluster to grow $\Rightarrow$ different $p_c$.}
\end{frame}

\demoframe{GIF}{run - \texttt{lattices\_merged.gif}}

\begin{frame}{Defining the critical parameter $p_c$}
  Fix a lattice. Let
  \[
    \theta(p) \;=\; \P_p\!\big(\, 0 \leftrightarrow \infty \,\big)
  \]
  be the probability that the cluster of the origin is infinite.
  $\theta$ is monotone non-decreasing in $p$, and
  \[
    p_c \;:=\; \sup\{\, p \in [0,1] \,:\, \theta(p) = 0 \,\}
        \;=\; \inf\{\, p \in [0,1] \,:\, \theta(p) > 0 \,\}.
  \]

  \begin{itemize}
    \item $p < p_c$: with probability $1$, every cluster is finite
          (sub-critical phase).
    \item $p > p_c$: a unique infinite cluster exists almost surely
          (super-critical phase).
    \item $p = p_c$: \emph{critical phase}. No infinite cluster, but
          clusters of every scale --- the scale-invariant regime where
          $\E V(r) \sim r^{d_f}$ lives.
  \end{itemize}

  \vspace{0.4em}
  \begin{center}
    \begin{tabular}{lccc}
      \toprule
      lattice              & honeycomb            & square                       & triangular \\
      vertex degree        & $3$                  & $4$                          & $6$        \\
      $p_c$                & $\approx 0.6962$     & $\approx 0.59274621$         & $= 0.5$    \\
      \bottomrule
    \end{tabular}
  \end{center}

  \TODO{Higher vertex degree $\Rightarrow$ easier for the cluster to
        grow $\Rightarrow$ lower $p_c$. Triangular $p_c = 1/2$ is
        \emph{exact} (Sykes--Essam duality); square is only known
        numerically. See~\cite{grimmett1999percolation}.}
\end{frame}

% =====================================================================
\section{Critical exponents, scaling, universality}
% target: ≈ 10 min
% =====================================================================

\begin{frame}{The zoo of critical exponents}
  \footnotesize
  \begin{center}
  \setlength{\tabcolsep}{4pt}
  \renewcommand{\arraystretch}{1.35}
  \begin{tabular}{c l c l c l}
    \toprule
    & \textbf{Function} & & \textbf{Behavior} & \textbf{Exp.} &
      \textbf{Scaling rel.} \\
    \midrule
    \rowcolor{InvisibleBlue}
    $\boldsymbol\Rightarrow$ & Cluster volume at $p_c$           & $V(r)$        & $\E V(r) \sim r^{d_f}$                                                  & $d_f$            & $d_f = d/(\tau - 1)$ \\
    & Cluster number density           & $n(s,p)$      & $n(s,p) \sim s^{-\tau}\, f\!\big(s\,|p-p_c|^{1/\sigma}\big)$             & $(\tau, \sigma)$ & --- \\
    & Percolation probability          & $\theta(p)$   & $\theta(p) \sim (p-p_c)^{\beta}$                                         & $\beta$          & $\beta = (\tau - 2)/\sigma$ \\
    \rowcolor{InvisibleBlue}
    $\boldsymbol\Rightarrow$ & Mean cluster size                & $\chi^f(p)$   & $\chi^f(p) \sim |p-p_c|^{-\gamma}$                                       & $\gamma$         & $\gamma = (3 - \tau)/\sigma$ \\
    & Number of clusters per vertex    & $\kappa(p)$   & $\kappa(p) \sim |p-p_c|^{2-\alpha}$                                      & $\alpha$         & $\alpha = 2 - (\tau-1)/\sigma$ \\
    & Cluster moments                  & $\chi_k^f(p)$ & $\chi_{k+1}^f(p)/\chi_k^f(p) \sim |p-p_c|^{-\Delta}$                     & $\Delta$         & $\Delta = 1/\sigma = \gamma + \beta$ \\
    \bottomrule
  \end{tabular}
  \end{center}

  \vspace{0.4em}
  \normalsize
  \TODO{Verbal: this is the landscape --- five (six) observables, five
        (six) exponents, all tied by scaling relations. We will work
        on the two highlighted rows: the on-critical cluster volume
        (the running example, exponent $d_f$) and, off-critical, the
        mean cluster size (exponent $\gamma$). They are the easiest
        case: \emph{positive} exponents of the $r^\beta$ family with
        $\beta > 0$. The same statistical machinery lifts to the rest
        with extra work. Hyperscaling $d_f = d/(\tau-1)$ to be
        verified before the talk. Source: Stauffer 1979
        \cite{bigCluster}.}
\end{frame}

\begin{frame}{Two anchor observables}
  (a) cluster volume $\E V(r) \sim r^{d_f}$ at $p = p_c$.

  \vspace{0.4em}
  (b) mean cluster size $\chi(p) \sim |p - p_c|^{-\gamma}$
  (off-critical).

  \vspace{0.8em}
  \TODO{Both kinds of observables are important, but we focus on
        observables of the first kind --- positive exponents are
        easier to measure.

        Insert image: \texttt{fig\_running\_example.png}.}
\end{frame}

\begin{frame}{Why log--log linearises a power law}
  Heuristic:
  \begin{align*}
    \E V(r) &\sim r^{d_f} \\
    \log \E V(r) &\approx d_f \log r + \mathrm{const}.
  \end{align*}

  \TODO{Physicists ``read the slope''. Display
        \texttt{fig\_running\_example.png}.}
\end{frame}

\begin{frame}{Universality across lattices}
  \TODO{same $d_f$, three lattices, parallel lines.
        Embed \texttt{fig\_universality.png}.}
\end{frame}

\demoframe{GIF}{run - \texttt{percolation\_meanfield.gif}}

\begin{frame}{(Optional bridge to statistics) The classical CLT}
  \TODO{$S_n \sim \sqrt{n}$ as a ``toy universality statement'':
        exponent $1/2$ independent of step distribution (provided
        $\Var < \infty$). Bridge from physicists' universality to
        statisticians' universality.}
\end{frame}

% =====================================================================
\section{From physics heuristic to statistical question}
% target: ≈ 8 min
% =====================================================================

\begin{frame}{The running example, revisited}
  \TODO{re-display \texttt{fig\_running\_example.png}. Familiar image
        $\Rightarrow$ lower cognitive load when we start adding the
        statistical machinery.}
\end{frame}

\begin{frame}{Big dots ``trick''}
  \TODO{Compare \texttt{fig\_running\_example.png} with
        \texttt{fig\_running\_example\_small\_dots.png}.
        Same fit, same data; only the markersize changed. With smaller
        dots the eye stops interpolating across fat blobs and the
        finite-size deviation at small $r$ becomes visible. Moral:
        ``looks linear'' is a perceptual claim, not a statistical one.}
\end{frame}

\begin{frame}{What physicists measure vs.\ what they want}
  \TODO{They want $\E Y_i$, they observe $\overline Y_i$.
        $\E Y_i = a_0\, i^{\gamma}\exp\!\big(\sum_{j=1}^{J} a_j i^{-j}
        + \phi_J(i)\,i^{-J-1}\big)$ — finite-size expansion.
        Today: focus on $J = 1$.}
\end{frame}

\begin{frame}{Bias--variance trade-off in one picture}
  \TODO{Large $r$: small finite-size bias, expensive $\Rightarrow$ few
        samples $\Rightarrow$ high variance. Small $r$: cheap but
        biased. The log--log estimator must pick its scales accordingly.}
\end{frame}

\begin{frame}{Chalkboard: log--log linearisation}
  \TODO{board derivation of
        $\log\E Y_{\rho^k} = \log a_0 + k\beta\log\rho
        + a_1 \rho^{-k} + \phi_1(\rho^k)\rho^{-2k}$.
        Slide just states the resulting expansion.}
\end{frame}

% =====================================================================
\section{The log--log estimator and its three-part split}
% target: ≈ 6 min
% =====================================================================

\begin{frame}{The OLS slope estimator}
  \TODO{$\hat\beta = \sum_{k=m_0+1}^{m_0+m} w_{k,m}\,
        \log\overline Y_{\rho^k}$
        with $w_{k,m} = \dfrac{12\big(k - m_0 - (m+1)/2\big)}{m(m^2-1)}$.
        $w$ is the OLS coefficient on $k$.}
\end{frame}

\begin{frame}{The three-part split}
  \TODO{$\log\overline Y_{\rho^k} = \log\E Y_{\rho^k}
        + \log\big(\overline Y_{\rho^k}/\E Y_{\rho^k}\big)$.
        Combine: $\hat\beta - \beta = R(m,m_0) + S_{n,m} + Q_{n,m}$.
        $R$ = deterministic finite-size bias; $S$ = linear stochastic
        (CLT-able); $Q$ = quadratic Taylor remainder.}
\end{frame}

\begin{frame}{Linear regression as a linear functional}
  \TODO{Visual: $\hat\beta$ is a linear combination of the
        $\log\overline Y_{\rho^k}$'s with weights $w_{k,m}$ that sum
        to zero and integrate against $k$ to one. This is why the CLT
        machinery is applicable.}
\end{frame}

% =====================================================================
\section{Main theorem: CLT for $\hat\gamma$}
% target: ≈ 10 min
% =====================================================================

\begin{frame}{Main theorem}
  \TODO{Under $(2+\delta)$-moment, bounded $\phi_j$, and variance
        convergence:
        $\sqrt{n m^3}\,(\hat\gamma - \gamma)\;\Rightarrow\;
        \mathcal N\!\Big(0,\, \tfrac{12\sigma_\infty^2}{\log^2\rho}\Big).$
        Rate $\sqrt{nm^3}$ combines $n$ trials and $m$ scales.}
\end{frame}

\begin{frame}{Proof sketch — decomposition}
  \TODO{Recall $\hat\beta - \beta = R + S + Q$. Treat each piece
        separately. Strategy: $R \to 0$ via choice of $m_0$ (§6);
        $S$ drives the CLT; $Q$ vanishes faster than the CLT rate.}
\end{frame}

\begin{frame}{Proof sketch — $S_{n,m}$ via Lindeberg--Feller}
  \TODO{$S_{n,m}$ is a triangular array of weighted sums. Verify the
        Lyapunov $(2+\delta)$-condition $\Rightarrow$ CLT. State the
        limiting variance $12\sigma_\infty^2/\log^2\rho$.}
\end{frame}

\begin{frame}{Proof sketch — $Q_{n,m}$ Taylor remainder}
  \TODO{$\log(\overline Y / \E Y) - (\overline Y / \E Y - 1)$ is
        quadratic on the good event $\{|\overline Y - \E Y| \le
        \tfrac12 \E Y\}$. Good-event probability $\to 1$ fast;
        bad-event contribution is $o(1/\sqrt{nm^3})$.}
\end{frame}

\begin{frame}{Proof sketch — $R(m,m_0)$ finite-size bias}
  \TODO{Deterministic; $R \asymp \rho^{-m_0}/m^2$ under $J=1$ expansion.
        Made negligible by choosing $m_0$ large enough — but not too
        large, otherwise $S$ blows up. Foreshadow §6.}
\end{frame}

% =====================================================================
\section{The role of $m_0$ and Wilson confidence intervals}
% target: ≈ 7 min
% =====================================================================

\begin{frame}{Why throw away the first $m_0$ scales}
  \TODO{The first scales are the most biased; the OLS weights $w_{k,m}$
        give them large lever arms. Dropping them kills the leading
        $\rho^{-k}$ term in the expansion.}
\end{frame}

\begin{frame}{Wilson-style confidence interval}
  \TODO{Four terms, each with its own rate:
        $|\hat\gamma - \gamma| \le
        \underbrace{C_1\rho^{-m_0}/m^2}_{\text{finite size}}
        + \underbrace{C_2/(nm)}_{\text{Jensen good event}}
        + \underbrace{C_3 \exp(-cn)}_{\text{bad event}}
        + \underbrace{z_{1-\alpha/2}\,\sqrt{12\sigma_\infty^2/(\log^2\rho \cdot nm^3)}}_{\text{Gaussian}}.$
        Each term decays at a different rate — the knob is $m_0$.}
\end{frame}

\begin{frame}{Empirical comparison of $m_0$ regimes}
  \TODO{Embed \texttt{fig\_regime\_sweep.png} — $\hat d_f \pm$ SE
        (band: $\pm$ std) for $m_0 = $ const / $\alpha m /
        \tfrac12\log_\rho(nm^3)$.}
\end{frame}

\begin{frame}{Error vs.\ simulation time}
  \TODO{Embed \texttt{fig\_error\_vs\_time.png}; legend slopes per
        regime. Which $m_0$ schedule gives the best $|\mathrm{bias}|$
        decay rate per second of compute?}
\end{frame}

\begin{frame}{The masterpiece plot — error vs.\ budget $B$}
  \TODO{Embed L4 plot (to be produced): budget $B \in \{1,2,4,8,16\}$,
        $m(B) = m_0 + \log_2 B$, three regimes (const, $\alpha=1/2$,
        log-$\rho$). The ``which knob and how should you turn it''
        slide.}
\end{frame}

% =====================================================================
\section{Lower bound and optimality}
% target: ≈ 4 min
% =====================================================================

\begin{frame}{Minimax lower bound}
  \TODO{$\inf_{\hat\gamma}\sup_\theta \E_\theta[\hat\gamma -
        \gamma(\theta)]^2 \gtrsim B^{-2\omega/(d + 2\Omega)},$
        where $\omega$ is the leading exponent and $\Omega$ the
        correction-to-scale exponent. Proof: KL between two product
        measures differing in $\gamma$, then Fano.}
\end{frame}

\begin{frame}{Upper bound matches the lower bound up to logs}
  \TODO{Plug $m \asymp \log_\rho(B)$ and $n \asymp B/m$ in our CLT
        rate $1/\sqrt{nm^3}$; recover the lower bound up to log
        factors. $\Rightarrow$ the log--log estimator is near-minimax.}
\end{frame}

% =====================================================================
\section{(Optional) Closed-form sanity checks}
% target: ≈ 3 min (cut first if running long)
% =====================================================================

\begin{frame}{Bethe lattice}
  \TODO{$\E Y_i$ has a closed form via Stirling; verify the expansion
        used in §3 reproduces it.}
\end{frame}

\begin{frame}{Simple random walk}
  \TODO{Return-to-origin observable; closes the loop back to the
        classical-CLT analogy from §2.}
\end{frame}

% =====================================================================
\section{Wrap-up}
% target: ≈ 2 min
% =====================================================================

\begin{frame}{Take-aways}
  \TODO{(1) Physicists' log--log plot has a rigorous statistical theory.
        (2) The estimator is asymptotically near-optimal.
        (3) $m_0$ is the key practical lever; $\alpha m$ with
            $\alpha \in (1/4, 1/2)$ is robust in practice.
        (4) Confidence intervals: four explicit terms, each with a
            different rate — pick $m_0$ to balance.}
\end{frame}

\begin{frame}[plain]
  \vfill
  \begin{center}
    {\Huge Thank you.}\\[1em]
    {\large Questions?}
  \end{center}
  \vfill
\end{frame}

% =====================================================================
%  Appendix (post-Q&A backup frames)
% =====================================================================
\appendix

\begin{frame}{Appendix: pool simulation pipeline (L0--L4)}
  \TODO{One-pager on the hierarchical simulation structure (L0 trial
        $\to$ L1 mean $\to$ L2 slope $\to$ L3 replicas $\to$ L4 budget).
        Reference \texttt{notes/simulation\_hierarchy.md}.}
\end{frame}

\begin{frame}{Appendix: numerical precision \& BFS speed-up}
  \TODO{Box-clipped BFS vs.\ full Bernoulli + union-find; $20\times$
        speed-up at $N=4096$. Reference \texttt{coding/bench\_l0.py}.}
\end{frame}

\end{document}
